% Springer LNCS template. Needs llncs.cls and splncs04.bst (see paper/README.md).
% The \todo{...} macros mark text to fill in; take numbers from docs/FINDINGS_LOG.md.
\documentclass[runningheads]{llncs}

\usepackage[T1]{fontenc}
\usepackage[utf8]{inputenc}
\usepackage{graphicx}
\usepackage{booktabs}
\usepackage{amsmath,amssymb}
\usepackage{xcolor}
\usepackage{siunitx}
\usepackage[hidelinks]{hyperref}
\usepackage{url}

% Reminder macro: red in drafts. Redefine as {} before submission.
\newcommand{\todo}[1]{\textcolor{red}{[TODO: #1]}}
\newcommand{\addr}[1]{\texttt{#1}}

\begin{document}

\title{Characterizing and Mitigating Address Misuse Risks\\in Cross-Chain EVM Ecosystems}
% Optional running title: \titlerunning{Address Misuse in Cross-Chain EVM}

\author{First Author\inst{1} \and
Second Author\inst{1} \and
Third Author\inst{1}}
\authorrunning{F. Author et al.}
\institute{University / Faculty, City, Country\\
\email{\{author1,author2,author3\}@example.edu}}

\maketitle

\begin{abstract}
% 150--250 words. Structure: context -> gap -> approach -> main result -> contributions.
Users who send a transaction to an address that is not the contract they
intended lose funds or leave them stranded. Shao et al.\ recently characterized
such \emph{address misuse} on Ethereum and BNB Smart Chain. We partially
replicate their study and extend it with a \emph{cross-chain consistency}
analysis: an address that holds a contract on one chain but is empty on the
chain the user transacts on. \todo{Key quantitative result, e.g.\ share of
misuse addresses explained by cross-chain misalignment and loss covered.}
We further propose a risk score $R(A,C)$ and a wallet-side warning that runs
before signing, with \todo{latency} median overhead.
\todo{One sentence on gold-label quality, e.g.\ Cohen's $\kappa$.}
\keywords{Address misuse \and Cross-chain \and EVM \and Smart contracts \and Wallet security \and EIP-7702.}
\end{abstract}

%------------------------------------------------------------------
\section{Introduction}
\label{sec:intro}
% Paragraph 1: the problem (sending to an address without code / the wrong contract).
% Paragraph 2: prior work (Shao et al.) and its limits: per-chain only, no cross-chain root cause, no wallet-side warning.
% Paragraph 3: our approach.
\todo{Motivation paragraph with a concrete incident.}

\medskip\noindent\textbf{Research questions.}
\begin{itemize}
\item \textbf{RQ1} How well can the address-misuse rules of Shao et al.\ be replicated on raw transaction data? \todo{replication gap}
\item \textbf{RQ2} What fraction of misuse is explained by cross-chain address misalignment, and what are the root causes (same deployer and nonce, factory with CREATE2)?
\item \textbf{RQ3} Can a risk score computed before signing separate misuse from benign transactions at a low false-positive rate and low latency?
\end{itemize}

\medskip\noindent\textbf{Contributions.}
\begin{itemize}
\item A partial replication of address-misuse detection on Ethereum and BSC \todo{numbers}.
\item A cross-chain consistency analysis over \todo{N} addresses on eight networks, with a taxonomy (Types~1--6) and a root-cause breakdown.
\item A manually audited gold set of \todo{N} labels (Cohen's $\kappa=\todo{x}$) and negative sets N1--N8.
\item A risk score $R(A,C)$ and a MetaMask Snap that warns before signing, evaluated on \todo{setting}.
\end{itemize}

%------------------------------------------------------------------
\section{Background and Threat Model}
\label{sec:background}
\subsection{Address misuse}
% Define CA misuse (NCD, PCD) and EOA misuse (KEY, ADDR, EIP-7702) following docs/CODEBOOK.md section 3.
\todo{Define NCD, PCD, KEY, ADDR and EIP-7702 delegation~\cite{eip7702}; use the notation of the original paper.}

\subsection{Taxonomy}
% Table of Types 1--6.
\begin{table}[t]
\caption{Taxonomy of address misuse used in this work.}
\label{tab:taxonomy}
\centering
\begin{tabular}{@{}lll@{}}
\toprule
Type & Name & Defining condition\\
\midrule
T1 & Call to empty address & NCD, no contract elsewhere before the tx\\
T2 & Black-hole contract & Code present, no withdrawal path\\
T3 & Token sent to its own contract & \todo{condition}\\
T4 & Cross-chain misalignment & NCD here, contract elsewhere deployed earlier\\
T5 & Exposed key / sweeper / malicious 7702 & KEY, swept ADDR, or malicious delegate\\
T6 & Address poisoning & Look-alike address, zero-value lure\\
\bottomrule
\end{tabular}
\end{table}

\subsection{Threat model and ethics}
% Passive analysis only; never use leaked keys; truncate personal addresses.
We analyze public chain data passively. We never use leaked private keys, and
we truncate personal addresses in the text. \todo{Responsible disclosure statement.}

%------------------------------------------------------------------
\section{Related Work}
\label{sec:related}
% Shao et al. (base work); address poisoning; sweeper bots and key leakage; EIP-7702 security; pre-sign transaction simulation and wallet warnings; bridge security.
Shao et al.~\cite{shao2026lost} \todo{summarize contribution and limits}.
\todo{Related work on address poisoning, leaked-key sweepers, EIP-7702 delegation attacks and pre-sign transaction simulation.}

%------------------------------------------------------------------
\section{Data and Methodology}
\label{sec:method}
\subsection{Data sources}
% Author dataset (Zenodo), BigQuery crypto_ethereum (ETH), Etherscan V2 / NodeReal (BSC first tx), RPC eth_getCode on 8 networks, Binplorer (BSC creation).
\begin{itemize}
\item Author dataset (Zenodo) as seed: \todo{counts per chain and subtype}.
\item Raw transactions: BigQuery \texttt{crypto\_ethereum} for Ethereum; \todo{BSC source}.
\item Cross-chain state: \texttt{eth\_getCode} on eight networks (Ethereum, BSC, Arbitrum, Optimism, Base, Polygon, Sepolia, BSC testnet).
\item Contract creation: Etherscan \texttt{getcontractcreation}, and Binplorer for BSC.
\end{itemize}

\subsection{Rule replication}
% Port of Algorithms 1--3; paper_compat flags for bugs B1/B2; pilot of 500 ETH addresses matched 500/500 in snapshot mode.
We port Algorithms~1--3 of \cite{shao2026lost} and report two modes: a
\emph{snapshot} that reproduces the original data cut and a \emph{current} mode.
\todo{Describe known deviations (selector 0xffffffff, sweep pairing) and results.}

\subsection{Cross-chain consistency matrix}
% Address x chain matrix: deployer, CREATE nonce, factory, NCD-XC/PCD-XC labels, redeployability, timing.
For each misuse address we build an $\text{address}\times\text{chain}$ matrix of
code state, deployer, CREATE nonce, creation kind and creation time. We label an
address \emph{NCD-XC} (resp.\ \emph{PCD-XC}) when it is empty (resp.\ holds
different code) on the victim chain but holds a contract on another chain, and
classify whether it can be \emph{redeployed} (claimable, burned, factory-bound,
unknown). \todo{Formalize with notation.}

\subsection{Gold labels and negative sets}
% Codebook, two annotators, kappa; N1--N8.
Two annotators label a stratified sample independently following a codebook
(Types T1--T6); we report Cohen's $\kappa$ on misuse and on type.
Negative sets N1--N8 are: \todo{list: random unflagged addresses, benign 7702 delegates, multi-chain infrastructure contracts, \ldots}.

%------------------------------------------------------------------
\section{Risk Score and Wallet-Side Warning}
\label{sec:design}
\subsection{Risk score $R(A,C)$}
% R(A,C) in [0,100]; features: code here, code elsewhere and timing, calldata, symbolic execution, delegate, history.
Given a destination address $A$ and the transaction context $C$ (chain,
calldata, value), the score is $R(A,C)\in[0,100]$ computed from
\todo{features and weights, e.g.\ code state, cross-chain consistency, delegate, intent}.

\subsection{Wallet-side warning}
% MetaMask Snap + risk engine; latency budget 150 ms.
The warning runs as a MetaMask Snap \todo{transaction-insight API} that queries a
risk engine before the user signs, with a latency target of
\SI{150}{\milli\second}. \todo{Architecture figure.}

%------------------------------------------------------------------
\section{Evaluation}
\label{sec:eval}
\subsection{Replication (RQ1)}
\todo{Table: ours vs.\ author numbers (addresses, misuse tx, loss), agreement.}

\subsection{Cross-chain analysis (RQ2)}
% Existing numbers are in docs/FINDINGS_LOG.md: NCD-XC share, deployed_before share, redeployability breakdown.
\todo{Table: per chain, NCD-XC and PCD-XC share of addresses and of loss; timing (deployed before / after); redeployability breakdown.}

\subsection{Detection quality (RQ3)}
\todo{Precision, recall, F1, FPR on the gold set; ROC for $R(A,C)$; ablation of the cross-chain feature.}

\subsection{Latency}
\todo{Median and p95 latency of the Snap and the engine.}

%------------------------------------------------------------------
\section{Discussion}
\label{sec:discussion}
\todo{Implications for wallets and explorers. Limitations: current-state code versus state at transaction time, snapshots, BSC data coverage, API truncation, label subjectivity. Generalization to other EVM chains.}

\section{Conclusion}
\label{sec:conclusion}
\todo{Summary of findings and future work.}

\subsubsection*{Acknowledgments.} \todo{Funding, if any.}

%------------------------------------------------------------------
\bibliographystyle{splncs04}
\bibliography{references}

\end{document}
